\setcounter{section}{0}
\hypertarget{invited-indigenous-engagement-panel}{}
\label{invited-page-indigenous-engagement-panel}
\section{Fostering Connections and Planning Research Projects with Indigenous Communities}
\textbf{Session:} \hyperlink{schedule-tuesday-afternoon-a-1502}{Tuesday -- Afternoon}\\
\textbf{Room:} A-1502 \textbf{Plenary session}\\
%\textbf{Theme:} \tetxbf{Plenary session}\\



\paragraph{Abstract}
{\itshape
Fostering Connections and Planning Research Projects with Indigenous Communities


This panel brings together Olivier Hernandez, Raina Irons, Melanie Demers, and Wilfred Buck to explore approaches to building meaningful relationships and co-developing research projects with Indigenous communities. Through personal reflections, panelists will introduce themselves and share their perspectives on Indigenous astronomy and its cultural, scientific, and relational significance.

The discussion will address key challenges in collaborative work between Indigenous and non-Indigenous partners, including barriers rooted in academic structures, trust-building, and differing knowledge systems. Panelists will highlight concrete examples of projects that center community engagement and Indigenous methodologies, examining how these initiatives were designed to benefit communities in reciprocal and respectful ways.

The panel will also consider the risks of both engaging and failing to engage with communities, as well as academic practices that may discourage Indigenous participation in research. Finally, it will reflect on the benefits of holistic approaches that integrate Indigenous knowledges, offering pathways toward more inclusive, ethical, and impactful research practices.

Panelists:
Olivier Hernandez (Planétarium de Montréal), Raina Irons, Melanie Demers,
Wilfred Buck
}

\hypertarget{invited-vincent-henault-brunet}{}
\label{invited-page-vincent-henault-brunet}
\section{Vincent Hénault-Brunet}
\textbf{Session:} \hyperlink{schedule-tuesday-morning-a-4502}{Tuesday -- Morning}\\
\textbf{Room:} A-4502\\
\textbf{Theme:} Stars and Stellar cluster formation\\

\paragraph{Biography}
{\sffamily
Dr. Vincent Hénault-Brunet is an Associate Professor in the Department of Astronomy and Physics at Saint Mary's University, in Halifax, where he is also the Director of the Burke-Gaffney Observatory. He was previously a Plaskett Fellow at the Herzberg Astronomy and Astrophysics Research Centre, an Excellence Initiative Fellow at Radboud University (NL) and a Research Fellow at the University of Surrey (UK). He completed his PhD at the Institute of Astronomy, University of Edinburgh (UK). His work focuses on the internal dynamics of star clusters, in particular old globular clusters. He uses a mix of kinematic data, dynamical models, and modern statistical techniques to probe the formation and evolution of these systems and the black holes they contain.
}

\paragraph{Abstract}
{\itshape
Rewinding the Dynamical Evolution of Globular Clusters and Their Black Hole Populations

Globular clusters (GCs) are among the oldest stellar systems in the Universe and may be the surviving counterparts of the compact, dense star-forming clumps now observed at high redshift with JWST, with initial densities reaching up to about a million solar masses per cubic parsec. Such initial conditions are expected to produce large populations of tightly packed stellar remnants in cluster cores, including black holes (BHs), which play a central role in cluster dynamical evolution and are now recognized as important contributors to gravitational-wave sources detected by LIGO-Virgo-KAGRA. In addition, dense GCs provide possible environments for the formation of intermediate-mass black holes (IMBHs) and exotic transients such as fast radio bursts. The presence of BHs in GCs also has a significant impact on the properties of stellar streams formed from dissolving clusters and used to probe dark matter. However, major uncertainties remain regarding the initial conditions of GCs, BH natal kicks, and the retention and demographics of their BH populations at the present day.

In this talk, I will present our recent work aimed at “rewinding” the dynamical evolution of Milky Way GCs to address open questions about the formation and evolution of these systems. Using a combination of multimass dynamical modelling and a fast cluster evolution code, we model GCs by fitting a wide range of observations (spatial distribution of stars, kinematics, pulsar timing) as boundary conditions for the long-term evolution of clusters and their BH populations. I will show evidence that GCs formed with a bottom-light initial stellar mass function (IMF), implying enhanced production of massive remnants compared to the commonly assumed IMF. I will then discuss constraints on present-day BH populations across a large sample of clusters, and how these results inform both their early densities and the magnitude of BH natal kicks. I will present new modelling of the orbits of fast-moving stars in Omega Centauri, supporting the presence of an IMBH and enabling predictions for future observational tests, as well as new stringent upper limits on the central dark mass of other massive GCs based on pulsar timing. Together, these results connect present-day observables to the early evolution of GCs and provide new insight into the formation and role of BHs in dense stellar systems.
}

\hypertarget{invited-laurence-perreault-levasseur}{}
\label{invited-page-laurence-perreault-levasseur}
\section{Laurence Perreault-Levasseur}
\textbf{Session:} \hyperlink{schedule-tuesday-morning-a-3561}{Tuesday -- Morning}\\
\textbf{Room:} A-3561\\
\textbf{Theme:} Cosmology and Early Universe\\

\paragraph{Biography}
{\sffamily
Laurence Perreault-Levasseur currently holds the Canada Research Chair in Artificial Intelligence and Computational Cosmology. She is an associate professor at the Université de Montréal and an Associate Member of Mila, where she conducts research in the development and application of machine learning methods to cosmology, with a special emphasis on strong gravitational lensing and other high-dimensional inverse problems. She is also a Visiting Scholar at the Flatiron Institute in New York City. Prior to that, she was a Flatiron research fellow at the Center for Computational Astrophysics in the Flatiron Institute and a KIPAC postdoctoral fellow at Stanford University. Laurence completed her PhD degree at the University of Cambridge, where she worked on applications of open effective field theory methods to the formalism of inflation. She received her B.Sc. and M.Sc. degrees from McGill University.
}

\paragraph{Abstract}
{\itshape
Cosmology from Learning the Universe: Results, Lessons, and Prospects

The Simons Collaboration on Learning the Universe (LtU) is a large, interdisciplinary effort that develops new computational and machine-learning–based methods to connect fundamental physics with observations of the universe. Its annual meeting brings together collaboration members and external experts to present advances in areas such as cosmological simulations, galaxy formation, and inference techniques, while fostering links with major survey projects.

More broadly, LtU aims to build fast, physically grounded models that bridge theory and data, enabling scientists to extract cosmological insights from next-generation surveys with greater efficiency than traditional simulations allow.
}

\hypertarget{invited-renee-hlozek}{}
\label{invited-page-renee-hlozek}
\section{Renée Hložek}
\textbf{Session:} \hyperlink{schedule-tuesday-afternoon-a-1502}{Tuesday -- Afternoon}\\
\textbf{Room:} A-1502\\
\textbf{Theme:} Next generation observatories - I\\

\paragraph{Biography}
{\sffamily
Renée Hložek is an Associate Professor at the David A. Dunlap Department of Astronomy \& Astrophysics and the Dunlap Institute at the University of Toronto. A leading expert in theoretical and observational cosmology, her research focuses on understanding the composition, structure, and evolution of the Universe through observations of the cosmic microwave background, dark matter, dark energy, and large astronomical surveys. She is a member of several major international collaborations, including the Atacama Cosmology Telescope, the Simons Observatory, and the Vera C. Rubin Observatory Dark Energy Science Collaboration, where she has played key leadership roles. Originally from South Africa, she completed her PhD at the University of Oxford as a Rhodes Scholar before joining Princeton University as a Lyman Spitzer Fellow and later moving to Toronto. Her work combines astrophysics, statistics, and machine learning to address some of the most fundamental questions in cosmology.
}

\paragraph{Abstract}
{\itshape
Next-generation science with the Rubin Observatory and CanDIAPL

The Vera C. Rubin Observatory has opened her massive eye on the cosmos, and in the coming months will begin the Legacy Survey of Space and Time (LSST). Rubin will revolutionize optical wide-field astronomy while entering a new era of all-sky short-period time-domain astronomy in the optical. Canadian scientists have access to Rubin data through in-kind contributions of a LSST Independent Data Access Center (IDAC), a public archive at the Canadian Astronomy Data Centre (CADC), and software infrastructure provided by the Canadian Rubin Fellows program, supported by the University of Waterloo and the Dunlap Institute for Astronomy and Astrophysics. To facilitate new science with this amazing data, we are developing the Canadian Data-Intensive Astrophysics PLatform (CanDIAPL, pronounced 'candy apple'), which is a platform integrated into the Canadian Advanced Network for Astronomical Research (CANFAR). CanDIAPL will co-host the Rubin data alongside other multi-wavelength data sets including radio data from MeerKAT and ASKAP, the precursor surveys to the Square Kilometer Array (SKA) and gravitational wave data sets from LIGO/Virgo/Kagra. The CanDIAPL team is developing software infrastructure that will enable cutting edge science with these data sets for all Canadians. I will describe some of the transformative science with the Rubin Observatory; with the start of the Rubin alert stream, some of that starts now as Rubin discovers new moving and changing objects in the sky, including asteroids, variable stars, supernovae, and active galactic nuclei! I will highlight some of the CanDIAPL tools being developed (that will also be demonstrated at our CanDIAPL booth at CASCA) and outline the steps for access to Rubin data for the Canadian community.
}

\hypertarget{invited-jason-hessels}{}
\label{invited-page-jason-hessels}
\section{Jason Hessels}
\textbf{Session:} \hyperlink{schedule-tuesday-afternoon-a-4502}{Tuesday -- Afternoon}\\
\textbf{Room:} A-4502\\
\textbf{Theme:} Transients\\

\paragraph{Biography}
{\sffamily
Jason Hessels is the Canada Excellence Research Chair (CERC) in Transients Astrophysics and Professor of Physics at McGill University and the Trottier Space Institute. He is also affiliated with the University of Amsterdam's Anton Pannekoek Institute for Astronomy and ASTRON, the Netherlands Institute for Radio Astronomy. He founded and co-leads the AstroFlash (https://astroflash-frb.github.io) research team, which focuses on studying the astrophysics of coherent radio transients, and using them to probe the nature of gravity, ultra-dense matter, and the otherwise invisible plasma and magnetic fields between stars and galaxies. He is using CHIME, CHORD, and LOFAR2.0 to open a new parameter space for transient discovery and many other telescopes worldwide to decipher the nature of these enigmatic sources --- including fast radio bursts (FRBs), long-period transients, and pulsars. This work is enabled by pushing astronomical observations to the highest-possible time and spatial resolution, while also casting a broad net with wide-field instruments.
}

\paragraph{Abstract}
{\itshape
Fast flashes from far-far away

In the past decade, we've learned that the sky is buzzing with fast radio bursts (FRBs) that last for mere milliseconds and originate from galaxies at millions to billions of parsecs. These remarkable bursts are trillions of times more luminous than Galactic pulsars, and yet the phenomenon is not rare. The wide-field CHIME radio telescope is discovering several FRB sources per day, and upcoming facilities like CHORD will discover tens of FRBs per day. We aim to understand the astrophysical origins of these enigmatic bursts and to use them as unique probes of the magneto-ionised media between stars and galaxies. While the ultra-magnetic neutron stars known as "magnetars" are likely responsible for at least some of the FRBs we see, the diversity of FRB environments and signal properties suggests that the observed phenomenon has multiple astrophysical origins. If so, that makes the astrophysical puzzle even deeper and richer. In this plenary talk, I will present our state-of-the-art understanding of FRBs and their applications as astrophysical and cosmological probes. I will also show that we've barely scratched the surface in terms of what we can find: there is likely an even larger population of coherent radio transients that we are on the verge of discovering with our increasingly powerful radio telescopes.
}

\hypertarget{invited-kim-venn}{}
\label{invited-page-kim-venn}
\section{Kim Venn}
\textbf{Session:} \hyperlink{schedule-tuesday-afternoon-a-3561}{Tuesday -- Afternoon}\\
\textbf{Room:} A-3561\\
\textbf{Theme:} Gas, dust and the ISM\\

\paragraph{Biography}
{\sffamily
Kim Venn is a Professor of Physics \& Astronomy at the University of Victoria. She was the founding Director of the UVic Astronomy Research Centre (2015-2024) to bring together scientists and engineers from UVic, NRC-Herzberg, and TRIUMF. During that time, she also led an NSERC CREATE program on New Technologies for Canadian Observatories that impacted over 70 Canadian students in astronomy, physics, and engineering. She was Board Chair for the Gemini Observatories (2024-2025) and ACURA (2022-2024). As co-PI for the Gemini GHOULS Large and Long program, she is studying the origin of the elements and dark matter properties using metal-poor stars in Local Group satellites, and is a strong supporter of the E-ELT ANDES spectrograph.
}

\paragraph{Abstract}
{\itshape
The surfaces of stars preserve a fossil record
of the chemical composition of their birth clouds,
and over time, these become powerful tracers of
galactic evolution and the origin of the elements.
Stars in the littlest dwarf galaxies are particularly
valuable because (1) their chemical abundances reflect
fewer previous star formation events, offering clearer
insights into nucleosynthetic processes; and (2) their
orbital dynamics trace the gravitational potential of
their host systems, providing key tests of dark matter
Properties. In this talk, I will discuss recent work to
constrain the nature of dark matter in these systems
and to advance our understanding of the origins of the
heavy elements. I will also highlight the impacts of
forthcoming instrumentation on this field, ranging
from exciting new spectroscopic surveys to studies
enabled by ELT/ANDES.
}


\hypertarget{invited-nathalie-ouellette}{}
\label{invited-page-nathalie-ouellette}
\section{Nathalie N.-Q. Ouellette}
\textbf{Session:} Tuesday -- Evening\\
\textbf{Room:} A-1502\\
\textbf{Theme:} Public Talk\\

\paragraph{Biography}
{\sffamily
Nathalie Nguyen-Quoc Ouellette is an astrophysicist, science communicator and lifetime lover of all things space! She obtained her Ph.D. in Physics and Astronomy at Queen's University in Kingston, Ontario in 2016. Her research focuses on galaxy formation and evolution, particularly those found in clusters. Nathalie is currently the Deputy Director of the Trottier Institute for Research on Exoplanets (iREx) and the Mont-Mégantic Observatory (OMM) at the University of Montréal and is also the Outreach Scientist for the James Webb Space Telescope in Canada collaborating with the Canadian Space Agency. She is a frequent contributor and analyst in Canadian media on everything related to space. She also organises and participates in science outreach events from local to international scales to encourage the interest and participation of youth and the general public in space science and to increase scientific literacy in Canada.
}

\paragraph{Abstract}
{\itshape
Les astronomes savent depuis un siècle que les galaxies, d'immenses regroupements d'étoiles et de poussière, peuplent l'Univers bien au-delà de notre propre Voie Lactée. Le télescope spatial Hubble en a révélé toute l'immensité grâce à ses images profondes, parsemées de milliers de galaxies. Aujourd'hui, le télescope spatial James Webb bouleverse à son tour notre compréhension de leur formation et de leur évolution, en dévoilant une surprenante surabondance de galaxies massives dans l'Univers primordial. Découvrez comment Webb sonde ces monstres célestes avec son regard infrarouge, et comment des astronomes, dont plusieurs au Canada, participent à redéfinir notre compréhension de la structure de l'Univers.
}


\hypertarget{invited-laurie-rousseau-nepton}{}
\label{invited-page-laurie-rousseau-nepton}
\section{Laurie Rousseau-Nepton}
\textbf{Session:} \hyperlink{schedule-wednesday-morning-a-4502}{Wednesday -- Morning}\\
\textbf{Room:} A-4502\\
\textbf{Theme:} Galaxy\\

\paragraph{Biography}
{\sffamily
Laurie Rousseau-Nepton is a faculty at the University of Toronto and the Dunlap Institute for Astronomy and Astrophysics. She comes with six years of experience working as a resident astronomer at the Canada-France-Hawaii Observatory supporting various instruments including wide-field cameras, high-resolution spectrographs, Fourier Transform Spectro-imager. She received her diploma from Université Laval by studying regions of star formation in spiral galaxies and helping with the development of two Fourier Transform Spectro-imagers, SpIOMM and SITELLE. She is now leading an international project called SIGNALS, the Star formation,Ionized Gas, and Nebular Abundances Legacy Survey, which sampled with the SITELLE instrument more than 50,000 of star-forming regions in 31 nearby galaxies to understand how the local environment affect the young star clusters characteristics.
}

\paragraph{Abstract}
{\itshape
Millions of Futures in the study of Galaxies
Understanding galaxy evolution requires connecting physical processes across a vast range of physical scales, from the internal dynamics of the interstellar medium to the large-scale environments in which galaxies reside. Recent discoveries have highlighted the importance of feedback processes, gas accretion, and environmental interactions in shaping galaxy morphology, star formation histories, and chemical enrichment. High-resolution observations and multi-wavelength surveys have revealed complex structures within galaxies, such as the complex morphologies of star-forming regions, outflows, and kinematic substructures, while also establishing statistical trends across galaxy populations over cosmic time. In this context, the study of galaxies at different scales has become increasingly integrated, combining spatially resolved spectroscopy, deep imaging, and simulations to build a coherent picture of galaxy assembly. The SIGNALS survey plays a key role in this effort by providing detailed, spatially resolved observations of nearby galaxies, enabling direct links between small-scale physical processes and global galaxy properties. Looking ahead, current and upcoming telescope and instrument capabilities, including next-generation integral field spectrographs and large-aperture observatories, will significantly enhance sensitivity, spatial resolution, and spectral coverage, opening new windows onto the multi-scale processes that drive galaxy evolution across cosmic history.
}

\hypertarget{invited-charles-edouard-boukare}{}
\label{invited-page-charles-edouard-boukare}
\section{Charles-Édouard Boukaré}
\textbf{Session:} \hyperlink{schedule-wednesday-morning-a-3561}{Wednesday -- Morning}\\
\textbf{Room:} A-3561\\
\textbf{Theme:} Exoplanets - I\\

\paragraph{Biography}
{\sffamily
Charles-Édouard Boukaré's research focuses on understanding the deep interior structure and evolution of rocky planets. His work lies at the intersection of earth and planetary sciences, computational fluid mechanics, mineral physics, and high-pressure chemistry.

He earned dual master's degrees in Engineering Geology and Planetary Sciences in 2012 from the École Nationale Supérieure de Géologie de Nancy (France) before completing a PhD in Geophysics at Université de Lyon (France) in 2016. He then pursued postdoctoral research at Brown University (USA), EPFL (Switzerland), IPGP (France), and McGill University (Canada). He has been an Assistant Professor in the Department of Physics and Astronomy at York University (Canada) since 2023.

Boukaré is particularly known for his work on magma oceans and early mantle dynamics. Over the past ten years, his innovative computational fluid dynamics models and thermodynamic analyses have provided new insights into the early evolution of the Earth, the Moon, Mercury, and exoplanets.

In 2022, he was awarded the prestigious international Doornbos Prize for outstanding contributions to the study of the Earth's deep interior. This honor, presented by the Committee on Studies of the Earth's Deep Interior (SEDI) of the International Union of Geophysics and Geodesy, recognizes exceptional work by early-career scientists.
}

\paragraph{Abstract}
{\itshape
Investigating the Interiors of Rocky Planets with the James Webb Space Telescope

Rapid advances in exoplanet research and planetary exploration are blurring the traditional boundaries between Earth sciences, geology, planetary science, and astronomy. Even within the Solar System, planetary interiors remain difficult to access, yet they play a fundamental role in shaping surface conditions. On Earth, for example, plate tectonics and the presence of a magnetic field are direct manifestations of internal dynamics. If probing planetary interiors is already challenging within the Solar System, how can we hope to constrain the interior properties of exoplanets?

With the launch of the James Webb Space Telescope, we have entered a new era in the characterization of exoplanets. In this presentation, I will provide a specific example of how combining geophysical concepts with state-of-the-art observational techniques can help reveal the interior state of rocky planets. I will focus in particular on a class of exoplanets known as lava worlds.
}

\hypertarget{invited-john-ruan}{}
\label{invited-page-john-ruan}
\section{John Ruan}
\textbf{Session:} \hyperlink{schedule-wednesday-afternoon-a-1502}{Wednesday -- Afternoon}\\
\textbf{Room:} A-1502\\
\textbf{Theme:} Compact Objects\\

\paragraph{Biography}
{\sffamily
John is an Associate Professor at Bishop's University in Sherbrooke, Quebec. He studies time-variable electromagnetic emission from gravitational wave sources, as well as their host galaxy properties. His research interests span across both the electromagnetic spectrum and the gravitational wave spectrum.
}

\paragraph{Abstract}
{\itshape
Gravitational wave (GW) experiments such as Pulsar Timing Arrays and the Laser Interferometer Space Antenna aim to detect supermassive black hole binaries at the centres of galaxies for the first time. Identifying the electromagnetic counterpart to each GW source is key to unlocking a variety of science goals for these experiments, although it is currently unclear how to approach this challenging task, primarily because the electromagnetic signatures of supermassive black hole binaries are still unclear. I will discuss the challenges to performing telescope follow-up for GW detections of supermassive black hole binaries, and some recent advances. Specifically, I will discuss (1) possible transient electromagnetic emission from the mergers of these binaries and the contamination from unrelated transients in the GW localization region, (2) detecting persistent periodic light curve variability from an actively-accreting binary and how to distinguish them from normal AGN, and (3) the distinct properties of their host galaxies in optical imaging and spectroscopy. Telescope follow-up of supermassive black hole binaries detected in gravitational waves will likely involve a combination of all of these approaches, and require expertise from a broad range of astronomers.
}

\hypertarget{invited-frederique-baron}{}
\label{invited-page-frederique-baron}
\section{Frédérique Baron}
\textbf{Session:} \hyperlink{schedule-wednesday-afternoon-a-3561}{Wednesday -- Afternoon}\\
\textbf{Room:} A-3561\\
\textbf{Theme:} Education and Public Outreach - I\\

\paragraph{Biography}
{\sffamily
Frédérique Baron is an astrophysicist, science communicator, and instrumentation project manager at the Laboratoire d'astrophysique expérimentale of the Observatoire du Mont-Mégantic. She also is the coordinator of the Centre for Research in Astrophysics of Québec (AstroQuébec) and has played a leading role in education and public outreach initiatives in Canadian astronomy. Her work bridges astrophysics research, instrumentation development, and science communication.
}

\paragraph{Abstract}
{\itshape
Canada is entering a pivotal period for astrophysics education and public outreach. New observatories, large international collaborations, and upcoming space missions are creating opportunities to rethink how Canadian astronomy connects with the public. This panel will discuss how outreach efforts are currently organized across Canada and what structures, partnerships, and priorities may be needed for the next decade.

The discussion will begin with an overview of the Canadian outreach landscape, including the role of initiatives such as the Westar program and national coordination efforts within the Canadian astronomical community. Panelists will discuss current initiatives, challenges in sustaining long-term programs, and the growing need for national coordination in education and public outreach (EPO).

The panel will also examine how EPO can be integrated into major Canadian-led astrophysics projects and infrastructures, such as the Very Large Optical Telescopes, future opportunities linked to upcoming space missions and observatories such as ARIEL, CASTOR, and the Habitable Worlds Observatory, as well as broader efforts to connect frontier astrophysics with classrooms and communities across Canada.

Finally, we will focus on the future: what should Canadian astrophysics outreach look like in 10 years, and how do we build sustainable, inclusive, and nationally connected programs that can support that vision?

Panelists:

Dan Morrone, Westar coordinator
Julie Bolduc-Duval, director of the Canadian astronomy training program Discover the Universe
Nathalie Nguyen-Quoc Ouellette, Outreach Scientist for the James Webb Space Telescope, Deputy Director of the Trottier Institute for Research on Exoplanets and the Observatoire du Mont-Mégantic
}

\hypertarget{invited-david-stenning}{}
\label{invited-page-david-stenning}
\section{David Stenning}
\textbf{Session:} \hyperlink{schedule-thursday-morning-a-1502}{Thursday -- Morning}\\
\textbf{Room:} A-1502\\
\textbf{Theme:} Astrostatistics, Astroinformatics, and AI / Machine Learning\\

\paragraph{Biography}
{\sffamily
Prof. Stenning is an Assistant Professor in the Department of Statistics \& Actuarial Science at Simon Fraser University (SFU). His main research area is astrostatistics, encompassing Bayesian methods, machine learning, and computer model emulation, with applications in exoplanetary science, stellar evolution, and solar physics, among others. He earned his PhD in Statistics from the University of California, Irvine in 2015 and held positions at the Institut d'Astrophysique de Paris, the Statistical and Applied Mathematical Sciences Institute via Duke University, and Imperial College London before joining SFU in January 2020.
}

\paragraph{Abstract}
{\itshape
Since Gauss's development of least squares in the early 1800s, challenges in astronomy have driven advances in statistics. This remains especially true today, as modern observing facilities produce large volumes of complex astrophysical data. These data are often noisy and high-dimensional, with heteroskedastic errors, selection effects, or missing observations, and they may require computationally expensive simulations for model comparison. Standard off-the-shelf methods often ignore these complexities, leading to biased inferences, poor predictive performance, and improper uncertainty quantification.

In this talk, I present several recent collaborative efforts in which astrophysical challenges motivated the development of new, general-purpose statistical methods. These include: (1) a machine learning approach for handling covariate shift between training and test data, applied to supernova classification and weak-lensing photometric redshift calibration; (2) a method for modeling noisy, nonhomogeneous Poisson processes, applied to identifying repeating fast radio bursts; and (3) a framework for emulating conditional distributions from stochastic simulations, applied to population studies of planetary system architectures.
}

\hypertarget{invited-kristine-spekkens}{}
\label{invited-page-kristine-spekkens}
\section{Kristine Spekkens}
\textbf{Session:} \hyperlink{schedule-thursday-morning-a-4502}{Thursday -- Morning}\\
\textbf{Room:} A-4502\\
\textbf{Theme:} Next generation observatories - II\\

\paragraph{Biography}
{\sffamily
Kristine Spekkens is a professor and Tier 1 Canada Research Chair in Gas-Rich Galaxy Structure in the Department of Physics, Engineering Physics and Astronomy at Queen's University. She leads a research group that focusses on understanding the structure and evolution of nearby galaxies, and she is particularly interested in using their atomic gas (HI) reservoirs as a cosmological probe. Spekkens leads a variety of programs to survey HI in galaxies using the biggest radio telescopes in the world, including facilities in Canada, Australia, the US and South Africa. Spekkens is also deeply involved in the international partnership to build the SKA Observatory (SKAO). Nationally, she is the Canadian SKA Science Director who focusses on coordinating across stakeholders to maximize the benefit of Canada's SKA investment. Internationally, she is the Chair of the SKA Science and Engineering Advisory Committee (SEAC), the SKAO Council's independent scientific and technical advisory body composed of specialists from around the world. Spekkens is also a strong advocate for equity, diversity and inclusion (EDI) in physics and astronomy, playing leading roles in national and international EDI initiatives.
}

\paragraph{Abstract}
{\itshape
Canada has a decades-long history of significant contributions to the global SKA initiative, culminating in our full membership in the SKA Observatory in 2024. With SKA construction well underway and the first science data expected next year, now is an exciting time to join the national and international SKA conversation. This talk will describe Canada and the SKA, emphasizing its tremendous scientific prospects using an example from my own research but also touching on the technology, digital research infrastructure, and community initiatives that are also fundamental to fulfilling the promise Canada's SKA investment.
}

\hypertarget{invited-marie-lou-gendron-marsolais}{}
\label{invited-page-marie-lou-gendron-marsolais}
\section{Marie-Lou Gendron-Marsolais}
\textbf{Session:} \hyperlink{schedule-thursday-afternoon-a-1502}{Thursday -- Afternoon}\\
\textbf{Room:} A-1502\\
\textbf{Theme:} Galaxy Clusters\\

\paragraph{Biography}
{\sffamily
Marie-Lou Gendron-Marsolais est astrophysicienne, professeure adjointe depuis 2024 au Département de physique, de génie physique et d'optique de l'Université Laval. Elle est titulaire de la chaire de recherche du Canada sur les environnements des amas de galaxies qui vise à décrire les influences internes (par le trou noir supermassif central) et externes (par l'environnement) sur les propriétés des galaxies pour améliorer notre compréhension de ces structures gravitationnellement liées. Pour se faire, elle analyse des observations multi-longueurs d'onde de pointe – radio, optique et rayons-X. Elle a obtenu son doctorat en astrophysique en 2018 à l'Université de Montréal. Elle a ensuite obtenu la prestigieuse bourse postdoctorale de l'Observatoire Européen Austral (ESO) qui lui a permis de poursuivre indépendamment ses propres travaux sur l'étude des amas de galaxies à Santiago, Chili, tout en travaillant pour le télescope ALMA (Atacama Large Millimeter/submillimeter Array). En 2021, elle s'est jointe en tant que chercheure postdoctorale à l'Instituto de Astrofísica de Andalucía, Grenade, Espagne, afin de coordonner les efforts des astronomes espagnols dans leur participation au projet du Square Kilometer Array Observatory (SKAO), qui sera le télescope radio le plus sensible jamais construit. Elle se passionne également pour la communication scientifique et les enjeux de la diversité en science.

Marie-Lou Gendron-Marsolais is an astronomer, assistant professor in the Department of Physics, Engineering Physics, and Optics at Université Laval since 2024. She holds the Canada Research Chair in Galaxy Cluster Environments, which aims to describe the internal (by the central supermassive black hole) and external (by the environment) influences on galaxy properties to improve our understanding of these gravitationally bound structures. To this end, she analyzes cutting-edge multi-wavelength observations—radio, optical, and X-ray. She obtained her PhD in astrophysics in 2018 from Université de Montréal. She then received the prestigious European Southern Observatory (ESO) postdoctoral fellowship, which allowed her to independently pursue her own research on galaxy clusters in Santiago, Chile, while working for ALMA (Atacama Large Millimeter/submillimeter Array). In 2021, she joined the Instituto de Astrofísica de Andalucía in Granada, Spain, as a postdoctoral researcher to coordinate the efforts of Spanish astronomers participating in the Square Kilometer Array Observatory (SKAO) project, which will be the most sensitive radio telescope ever built. She is also passionate about science outreach and solving the issues of diversity in science.
}

\paragraph{Abstract}
{\itshape
A radio eye on the Perseus cluster of galaxies and future prospects in SKAO era

Galaxy clusters are found at the intersections of the filaments of the Cosmic Web and form stormy, interacting environments, involving physics beyond the reach of a terrestrial laboratory. The advent of high-sensitivity low radio frequency facilities has recently shed new light on our conception of these systems, unveiling diffuse structures extending on large distances without a direct association with a host galaxy, demonstrating the prevalence of large-scale magnetic fields in these environments. Such sources have even been detected in a few clusters at higher redshifts, suggesting that magnetic fields are rapidly amplified in early cluster formation to offset inverse Compton losses. Closer to us, the notorious Perseus cluster is an astonishing case study for how high-quality images at low radio-frequencies can bring a whole new dimension to our understanding of a galaxy cluster. Indeed, VLA and LOFAR observations have revealed filamentary sub-structures in the radio lobes of the brightest cluster galaxy NGC 1275, the irregular morphology of its mini-halo and its connection with the bent-jet radio galaxy NGC 1272, the only known tailed radio galaxy with a gamma-ray counterpart (IC310), and the presence of a 1.1 Mpc large giant radio halo. I will review different interpretations of these structures, involving interaction with bulk motions of intracluster gas, galaxy's orbit in the cluster and Faraday depth measurements to reconstruct 3D structures. Such work is paving the way for the future generation of radio observatories, such as the world's next most sensitive radio telescope, the SKAO, which will uncover many more radio sources in these environments and open a new window on the unknown radio universe.
}

\hypertarget{invited-samantha-lawler}{}
\label{invited-page-samantha-lawler}
\section{Samantha Lawler}
\textbf{Session:} \hyperlink{schedule-thursday-afternoon-a-4502}{Thursday -- Afternoon}\\
\textbf{Room:} A-4502\\
\textbf{Theme:} Night-Sky Protection / Solar System\\

\paragraph{Biography}
{\sffamily
Samantha Lawler is an associate professor at the University of Regina. Her discoveries in the Kuiper Belt and predictions for satellite pollution have been featured by BBC, CBC, CNN, NPR, Scientific American, The New York Times, The Los Angeles Times, Wired Magazine, Nature, and many other international news outlets. She lives on a farm outside Regina and deeply appreciates the beautiful prairie skies.
}

\paragraph{Abstract}
{\itshape
The night sky is now changing dramatically thanks to the thousands of commercial satellites launched in the last several years. Because of these changes in the sky, we astronomers are having to modify how we do our research, how we stargaze, how we teach and learn astronomy, and how we reach out to the general public. Satellite pollution is adding to our workload with new and weird extra responsibilities that are being dumped on astronomers, such as:

responding to UFO sightings (thanks, Starlink trains, launches, and reentries)
tracking and identification of space debris on the ground in Canada (thanks, SpaceX)
learning and teaching international space law (thanks, SpaceX)
collisional cascades: not just for debris disks anymore (thanks, Starlink)
learning how to remove or avoid bright streaks in your data (thanks, Starlink)
learning how to avoid satellite RFI (thanks, Starlink direct-to-cell)
I'll talk about upcoming challenges, examples from planetary science research, and how we as a community can fight for dark and quiet skies above Canada.
}

